\input{../tex-inputs/latex-leseansicht-vorspann}

\section[Arthur und Olga Schnitzler an Gustav Schwarzkopf, 14. 11. 1906]{L04390 Arthur und Olga Schnitzler an Gustav Schwarzkopf, 14. 11. 1906}
\nopagebreak\mylabel{L04390v}
\rehead{ }\normalsize\beginnumbering\briefempfaengerindex{Schwarzkopf, Gustav@\textsc{Schwarzkopf, Gustav}!zzzSchnitzler, Olga@\emph{von Olga Schnitzler}!1906-11-142@{14. 11. 1906}|(be}\briefempfaengerindex{Schwarzkopf, Gustav@\textsc{Schwarzkopf, Gustav}!zzzSchnitzler, Arthur@\emph{von Arthur Schnitzler}!1906-11-142@{14. 11. 1906}|(be}
\toendnotes[C]{\smallbreak\pagebreak[2]}
\correspDesc{Versand  durch Arthur Schnitzler, Olga Schnitzler am 14. 11. 1906 in Semmering
\newline{}Übermittlung  am 14. 11. 1906 in Semmering
\newline{}Zustellung  am 15. 11. 1906 in 1/1 Wien 1
\newline{}Erhalt  durch Gustav Schwarzkopf am 15. 11. 1906 in Tiefer Graben 23}\toendnotes[C]{\smallbreak}
\Standort{DLA, A:Schnitzler, HS.1985.1.1897.}
\physDesc{Bildpostkarte, 113 Zeichen
\newline{}Handschrift Arthur Schnitzler: Bleistift, deutsche Kurrent
\newline{}Handschrift Olga Schnitzler: Bleistift
\newline{}Versand: 1) Stempel: »\nobreak{}\oindex{Semmering@\textbf{Semmering}|pwk}Semmering, 14. X{[}I. 1906{]}, 6\nobreak{}«.   2) Stempel: »\nobreak{}\oindex{I., Innere Stadt@\textbf{I., Innere Stadt}|pwk}1/1 Wien 1, 15. 11. 06, VIII\nobreak{}«. }\pstart{}{\pb}\textsc{Herrn Gustav Schwarzkopf}\pend{}\pstart{}\textsc{Wien I}\oindex{I., Innere Stadt@\textbf{I., Innere Stadt}|pw}\pend{}\pstart{}\textsc{Tiefer Graben 23}\oindex{Wien@\textbf{Wien}!I., Innere Stadt@\textbf{I., Innere Stadt}!Tiefer Graben 23@\textbf{Tiefer Graben 23}|pw}.\pend{}{\bigskip}
\pstart
           \centering{}{\pb}\textcolor{gray}{\textbf{Semmering, Sonnwendsteingipfel\oindex{Sonnwendstein@\textbf{Sonnwendstein}|pw} mit Schüler-Alpenhaus\oindex{Friedrich Schüler-Alpenhaus@\textbf{Friedrich Schüler-Alpenhaus}|pw}.}}\pend
           \vspace{1em}
\pstart
           \noindent{}{\pb}Herzliche Grüße.\pend
           
\pstart
           Sind hier\oindex{Semmering@\textbf{Semmering}|pw} mit  Brahm\pwindex{Brahm, Otto 5.\,2.\,1856 Hamburg – 28.\,11.\,1912 Berlin@\textsc{Brahm, Otto} (5.\,2.\,1856 Hamburg – 28.\,11.\,1912 Berlin)|pw} zuſammen.\pend
           
\pstart
           Ihr{\\[\baselineskip]}\spacefill\mbox{A. S.}\pend
           \leftskip=0em{}
\pstart
           14. 11. 906\pend
           \selectlanguage{ngerman}\vspace{1em}\pstart \spacefill\mbox{{[}hs. O. Schnitzler:{]} O. S.}\pend{}\selectlanguage{ngerman}\endnumbering\briefempfaengerindex{Schwarzkopf, Gustav@\textsc{Schwarzkopf, Gustav}!zzzSchnitzler, Olga@\emph{von Olga Schnitzler}!1906-11-142@{14. 11. 1906}|)be}\briefempfaengerindex{Schwarzkopf, Gustav@\textsc{Schwarzkopf, Gustav}!zzzSchnitzler, Arthur@\emph{von Arthur Schnitzler}!1906-11-142@{14. 11. 1906}|)be}\mylabel{L04390h}
\begin{anhang}
\end{anhang}\newcommand{\dateiname}{L04390}\newcommand{\titel}{Arthur und Olga Schnitzler an Gustav Schwarzkopf, 14. 11. 1906}\newcommand{\editorInnen}{Herausgegeben von Jahnke, SelmaMüller, Martin Anton}\input{../tex-inputs/latex-leseansicht-abspann}
